%*******************************************************
% 4.2 Tổng quan tài liệu
%*******************************************************
\section{Tổng quan những tài liệu liên quan đến tình hình nghiên cứu}\label{sec:tongquan}

\begin{tomtat}
Tài liệu được chia thành năm nhóm. Nhóm (A) giải thích vì sao tài chính tụ về một vài thành phố và các thành phố trong một nước phân vai ra sao. Nhóm (B) cho các thước đo thanh khoản, hiệu quả thông tin và minh bạch. Nhóm (C) cho bằng chứng khi cổ phiếu đổi sàn hoặc khi hai sàn hợp nhất (Tel Aviv, Euronext, Tokyo--Osaka). Nhóm (D) cho bộ công cụ danh mục khi tham số phải ước lượng. Nhóm (E) là nghiên cứu về Việt Nam, gần như chỉ về HOSE. Từ năm nhóm rút ra bốn khoảng trống \KK{1}--\KK{4} (Bảng~\ref{tab:khoangtrong}), và cả bốn cùng chỉ về một chỗ: chưa ai đo vai trò của sàn đối với chất lượng thị trường và khả năng đầu tư ở Việt Nam.
\end{tomtat}

Năm nhóm tài liệu dưới đây không được chọn theo chủ đề cho đủ, mà theo nhu cầu của các giả thuyết. Nhóm (A) phục vụ \HH{1} và phần định hướng cho Hà Nội; nhóm (B) phục vụ \HH{2} và thước đo của \HH{3}; nhóm (C) phục vụ \HH{3}; nhóm (D) phục vụ \HH{4}; nhóm (E) xác định đề tài đứng ở đâu so với các nghiên cứu về Việt Nam.

\subsection*{(A) Trung tâm tài chính: tích tụ, cạnh tranh và phân vai}

Kindleberger (1974) và Gehrig (2000) giải thích sự tập trung của hoạt động tài chính bằng lợi thế quy mô và thông tin. Lập luận cốt lõi là thông tin tài chính có giá trị nhất khi ở gần nơi nó được tạo ra và gần những người dùng nó, nên ngân hàng, quỹ và sàn giao dịch có lợi khi đặt cạnh nhau. Lợi thế này tự củng cố: trung tâm càng lớn thì càng hút thêm tổ chức tài chính.

Yao (2012) xác định các yếu tố quyết định sức hấp dẫn của trung tâm đối với ngành quỹ đầu tư qua trường hợp Luxembourg. Moosa và cộng sự (2016) kiểm tra độ vững của các yếu tố quyết định vị thế trung tâm, còn Eichengreen và Shah (2020) xếp các yếu tố thành năm nhóm, trong đó có nhóm phát triển tài chính và vốn hóa thị trường.

Naeem (2023) hồi quy thứ hạng GFCI của 196 trung tâm giai đoạn 2007--2020 theo 17 yếu tố chọn bằng phương pháp cây quyết định, không gồm chỉ báo thị trường cổ phiếu; Cetorelli và Peristiani (2013) đo vị thế trung tâm qua mạng lưới niêm yết quốc tế của sàn. Ở Trung Quốc, Lai (2012) và Wang và cộng sự (2007) mô tả sự phân vai giữa Bắc Kinh, Thượng Hải và Hồng Kông, Chen và Chen (2015) nhấn mạnh vai trò dẫn dắt của nhà nước, và Elliott (2011) cho rằng không có công thức chung để xây dựng một trung tâm toàn cầu.

Đối với đề tài, nhóm (A) cho hai thứ. Thứ nhất là thiết kế của \HH{1}: Naeem (2023) đã dựng sẵn khung dữ liệu bảng GFCI, nhưng bỏ ngoài chỉ báo thị trường cổ phiếu, nên việc thêm chỉ báo này vào là một bổ sung có thể kiểm tra được. Thứ hai là khung so sánh cho Hà Nội: trường hợp Bắc Kinh, Thượng Hải và Hồng Kông cho thấy các trung tâm trong một nước có thể chia vai thay vì cạnh tranh trên mọi mảng.

\subsection*{(B) Chất lượng thị trường vốn}

Thanh khoản thường được đo từ dữ liệu ngày bằng độ phi thanh khoản (Amihud, 2002), tỷ lệ ngày có lợi suất bằng 0 (Bekaert và cộng sự, 2007; Lesmond và cộng sự, 1999) và chênh lệch mua-bán ước lượng từ giá cao nhất và thấp nhất (Corwin \& Schultz, 2012); Lesmond (2005) và Fong và cộng sự (2017) đánh giá chỉ báo nào đo thanh khoản tốt nhất.

Ba thước đo nắm bắt ba mặt khác nhau của \enquote{dễ mua bán}. Độ phi thanh khoản Amihud trả lời câu hỏi: một đồng giá trị giao dịch làm giá dịch chuyển bao nhiêu. Tỷ lệ ngày lợi suất bằng 0 trả lời câu hỏi: có bao nhiêu ngày giá không nhúc nhích, thường vì không ai giao dịch hoặc vì chi phí giao dịch lớn hơn lợi ích của việc giao dịch. Chênh lệch Corwin--Schultz ước lượng khoảng cách giữa giá mua và giá bán từ giá cao nhất và thấp nhất trong ngày. Cả ba đều tính được từ dữ liệu ngày, nên áp dụng được cho cả giai đoạn dài từ năm 2009, kể cả khi dữ liệu trong ngày cho giai đoạn đầu không đầy đủ.

Hiệu quả thông tin được đo bằng tỷ số phương sai (Lo \& MacKinlay, 1988) và độ trễ giá (Hou \& Moskowitz, 2005), dù các thước đo này có thể sai lệch khi môi trường thông tin khác nhau (Griffin và cộng sự, 2010). Lang và cộng sự (2012) liên hệ minh bạch với thanh khoản, và Čihák và cộng sự (2012) xây dựng Cơ sở dữ liệu Phát triển Tài chính Toàn cầu. Các thước đo này chưa được quy đổi thành tổn thất của nhà đầu tư đa dạng hóa.

Câu cuối là chỗ đề tài bắt đầu. Biết một sàn có độ phi thanh khoản cao hơn sàn kia 30\% chưa cho biết nhà đầu tư mất bao nhiêu tiền. Đề tài quy đổi sự khác biệt về chất lượng thị trường thành chênh lệch lợi suất tương đương chắc chắn, một đơn vị mà nhà đầu tư và cơ quan quản lý đều hiểu được (mục~\ref{sec:khanangdautu}).

\subsection*{(C) Cấu trúc sàn, đổi sàn và hợp nhất sàn}

Huang và Stoll (1996) so sánh chi phí giao dịch của các cặp cổ phiếu tương đồng trên hai sàn, và Christie và Huang (1994) đo thanh khoản của cùng một cổ phiếu khi doanh nghiệp đổi sàn. Hai cách tiếp cận này là hai nửa của thiết kế nghiên cứu ở \HH{3}: so sánh cặp cổ phiếu tương đồng (nhóm đối chứng ghép cặp) và so sánh cùng một cổ phiếu trước và sau khi đổi sàn.

\paragraph{Tel Aviv.} Khi Sở Giao dịch Chứng khoán Tel Aviv chuyển cổ phiếu sang cơ chế có thêm phiên giao dịch liên tục, lợi suất bất thường tích lũy điều chỉnh theo thị trường đạt khoảng 5,5\% từ năm ngày trước công bố đến 30 ngày sau khi giao dịch theo cơ chế mới, và thanh khoản tăng cùng giá (Amihud và cộng sự, 1997). Đây là bằng chứng rằng thay đổi quy tắc giao dịch, không đổi doanh nghiệp, đủ làm thay đổi cả thanh khoản lẫn giá. Bài này còn ghi nhận thanh khoản lan sang cổ phiếu chưa chuyển cơ chế, một hiện tượng đề tài phải tính đến khi chọn nhóm đối chứng (mục~\ref{sec:did}).

\paragraph{Euronext.} Việc hình thành Euronext làm thanh khoản tăng ở doanh nghiệp lớn và doanh nghiệp có doanh thu nước ngoài, nhưng không tăng ở doanh nghiệp nhỏ (Nielsson, 2009). Bài học là tác động của hợp nhất không đồng đều: lợi ích dồn vào doanh nghiệp vốn đã được nhà đầu tư rộng rãi biết đến. Nếu điều này đúng ở Việt Nam, tác động bình quân của việc chuyển sàn sẽ phụ thuộc vào cơ cấu quy mô của các mã chuyển sàn, và mã nhỏ có thể được lợi ít; đề tài vì vậy ước lượng tương tác với quy mô.

\paragraph{Tokyo--Osaka.} Sau khi thị trường cổ phiếu của Osaka được hợp nhất vào Tokyo ngày 16/7/2013, thanh khoản của các cổ phiếu sẵn có trên Sở Giao dịch Chứng khoán Tokyo tăng trong ngắn hạn rồi giảm sau một năm, với hiệu ứng tâm lý mạnh hơn ở cổ phiếu nhỏ (Hisada, 2018); Arnold và cộng sự (1999, dẫn theo Hisada, 2018) cũng thấy hợp nhất sàn làm thanh khoản tăng. Trường hợp này gần với Việt Nam hơn Euronext ở chỗ là hợp nhất hai sàn trong một quốc gia. Cần chú ý Hisada đo cổ phiếu của sàn tiếp nhận (Tokyo), không đo cổ phiếu chuyển sang từ Osaka; đề tài đo nhóm ngược lại, tức các cổ phiếu chuyển sàn. Kết quả \enquote{tăng rồi giảm} cũng là lý do đề tài theo dõi đủ 24 tháng sau sự kiện thay vì vài tháng.

Wójcik (2007) phân tích quan hệ hai chiều giữa sàn chứng khoán và trung tâm tài chính. Bằng chứng ở thị trường mới nổi còn ít và chưa đo khả năng đầu tư.

\subsection*{(D) Đa dạng hóa danh mục khi tham số được ước lượng}

Tối ưu hóa trung bình-phương sai (Markowitz, 1952) dồn trọng số vào tài sản có sai số ước lượng lớn (Michaud, 1989), và hiệu quả ngoài mẫu giảm khi dùng ước lượng mẫu thay tham số tổng thể (Kan \& Zhou, 2007). Các cách khắc phục gồm co rút ma trận hiệp phương sai (Ledoit \& Wolf, 2004), mô hình Black-Litterman (Bessler và cộng sự, 2017; Black \& Litterman, 1992) và danh mục vững chắc (Garlappi và cộng sự, 2007).

Dù vậy, không mô hình nào vượt quy tắc 1/N một cách nhất quán (DeMiguel và cộng sự, 2009), và 1/N xấp xỉ tối ưu khi mức mơ hồ về mô hình cao (Pflug và cộng sự, 2012); Ledoit và Wolf (2008) kiểm định chênh lệch tỷ số Sharpe giữa hai danh mục. Đề tài dùng nhóm công cụ này để đo khả năng đầu tư của một sàn, không để chọn mô hình tốt nhất.

Nhóm (D) liên quan đến sàn giao dịch qua một con đường ít được chú ý. Mọi mô hình danh mục đều cần ước lượng lợi suất kỳ vọng và hiệp phương sai từ dữ liệu giá. Nếu cổ phiếu ít giao dịch, giá đứng yên nhiều ngày, thì ước lượng hiệp phương sai bị méo và sai số ước lượng lớn hơn. Một sàn có chất lượng thấp vì vậy không chỉ làm giao dịch đắt hơn, mà còn làm nhà đầu tư ước lượng kém hơn. Đề tài đo tổng của hai thiệt hại này qua lợi suất tương đương chắc chắn ngoài mẫu.

\subsection*{(E) Nghiên cứu về Việt Nam}

Vo và Phan (2017) tìm thấy hành vi bầy đàn trên các công ty niêm yết tại HOSE giai đoạn 2005--2015, còn Hoang và Luu (2024), trong một bản thảo công bố trực tuyến, phân tích chế độ biến động của VN-Index. Về danh mục, Ta và Vuong (2020), Nguyen và cộng sự (2020), Goutte và cộng sự (2025) và bài báo của chính thí sinh (Nguyễn \& Nguyễn, 2026a) áp dụng các mô hình Black-Litterman, co rút hiệp phương sai và ra quyết định đa tiêu chí trên cổ phiếu HOSE. Các công trình này chưa so sánh hai sàn.

\subsection*{Khoảng trống nghiên cứu}

Bảng~\ref{tab:khoangtrong} đối chiếu bốn nhóm công trình gần nhất với đề tài và nêu khoảng trống của từng nhóm; khoảng trống \KK{2} bổ sung cho ba vấn đề nêu ở mục~\ref{sec:lydo}.

\begin{table}[H]
\myfloatalign
\caption{Các nhóm nghiên cứu gần nhất và khoảng trống nghiên cứu}\label{tab:khoangtrong}
{\bangnho
\begin{tabularx}{\textwidth}{@{}>{\raggedright\arraybackslash}p{0.24\textwidth} L L@{}}
\toprule
Nhóm nghiên cứu & Kết quả đã có & Khoảng trống\\
\midrule
Trung tâm tài chính (Lai, 2012; Naeem, 2023; Yao, 2012) & Yếu tố quyết định thứ hạng GFCI (Naeem: 17 yếu tố, không có chỉ báo thị trường cổ phiếu); phân vai các trung tâm trong một quốc gia & \KK{1}. Các nghiên cứu về GFCI chưa xét chất lượng và cấu trúc thị trường cổ phiếu\\
\addlinespace
Chất lượng thị trường (Amihud, 2002; Fong và cộng sự, 2017; Lesmond, 2005) & Thước đo thanh khoản, hiệu quả thông tin, minh bạch ở thị trường mới nổi & \KK{2}. Chưa quy đổi chất lượng thị trường thành tổn thất của nhà đầu tư đa dạng hóa\\
\addlinespace
Đổi sàn và hợp nhất sàn (Hisada, 2018; Huang \& Stoll, 1996; Nielsson, 2009) & Đổi sàn làm thay đổi thanh khoản; Euronext: lợi ích dồn vào doanh nghiệp lớn; Tokyo--Osaka (cổ phiếu của sàn tiếp nhận): hiệu ứng tâm lý mạnh hơn ở cổ phiếu nhỏ & \KK{3}. Bằng chứng về hợp nhất sàn trong một quốc gia ở thị trường mới nổi còn ít, kết quả trái chiều theo quy mô, chưa đo khả năng đầu tư\\
\addlinespace
Việt Nam (Goutte và cộng sự, 2025; Nguyễn \& Nguyễn, 2026a; Vo \& Phan, 2017) & Nghiên cứu tập trung vào HOSE và VN-Index & \KK{4}. Chưa so sánh HNX với HOSE; chưa đo tác động của đợt chuyển sàn năm 2026\\
\bottomrule
\end{tabularx}}
\nguon{Nguồn: Thí sinh tổng hợp. Ghi chú: \KK{1}--\KK{4} là các khoảng trống, đánh số riêng để phân biệt với các nhóm tài liệu (A)--(E); phạm vi tìm kiếm nêu ở đoạn ngay dưới bảng.}
\end{table}

Bốn khoảng trống cùng chỉ ra một điểm: vai trò của sàn giao dịch đối với chất lượng thị trường và khả năng đầu tư chưa được đo cho Việt Nam, trong khi đợt hợp nhất năm 2026 tạo điều kiện đo trực tiếp. Nhận định này dựa trên tìm kiếm trên SSRN, RePEc/IDEAS, ScienceDirect, Oxford Academic, Wiley, Springer, Taylor \& Francis, MDPI và arXiv đến ngày 30/9/2026, chưa gồm Scopus, Web of Science và các tạp chí trong nước. Thí sinh mới đọc tóm tắt và thông tin thư mục của phần lớn các nguồn; việc đọc toàn văn và mở rộng tìm kiếm nằm trong kế hoạch năm thứ nhất.

Mỗi khoảng trống ứng với một giả thuyết: \KK{1} với \HH{1}, \KK{4} với \HH{2}, \KK{3} với \HH{3}, \KK{2} với \HH{4}. Cách ghép này giúp kiểm tra nhanh rằng đề cương không có giả thuyết nào thiếu lý do và không có khoảng trống nào bị bỏ rơi.

\begin{delam}
\emph{Dễ nghĩ là} câu \enquote{chưa có nghiên cứu nào} là một khẳng định tuyệt đối. Trong đề cương, khẳng định này luôn đi kèm phạm vi tìm kiếm (cơ sở dữ liệu nào, đến ngày nào) và giới hạn (mới đọc tóm tắt phần lớn nguồn). Khi trả lời hội đồng, nên nói \enquote{trong phạm vi tìm kiếm đến ngày 30/9/2026}, và sẵn sàng nhận rằng tìm kiếm trên Scopus, Web of Science và tạp chí trong nước là việc của năm thứ nhất.
\end{delam}
